\documentclass[10pt,a4paper]{amsart}
\usepackage{amsmath,amssymb,amsthm,graphicx}
\usepackage{verbatim}
\newenvironment{reference}{\comment}{\endcomment}
%\newenvironment{reference}{}{}
\newenvironment{slogan}{\comment}{\endcomment}
\newenvironment{history}{\comment}{\endcomment}

% For commutative diagrams we use Xy-pic
\usepackage[all]{xy}

% We use 2cell for 2-commutative diagrams.
\xyoption{2cell}
% Theorem environments.
%
\theoremstyle{plain}
\newtheorem{theorem}[subsection]{Theorem}
\newtheorem{proposition}[subsection]{Proposition}
\newtheorem{lemma}[subsection]{Lemma}

\theoremstyle{definition}
\newtheorem{definition}[subsection]{Definition}
\newtheorem{example}[subsection]{Example}
\newtheorem{exercise}[subsection]{Exercise}
\newtheorem{situation}[subsection]{Situation}

\theoremstyle{remark}
\newtheorem{remark}[subsection]{Remark}
\newtheorem{remarks}[subsection]{Remarks}

\numberwithin{equation}{subsection}

% Macros
%
\def\lim{\mathop{\mathrm{lim}}\nolimits}
\def\colim{\mathop{\mathrm{colim}}\nolimits}
\def\Spec{\mathop{\mathrm{Spec}}}
\def\Hom{\mathop{\mathrm{Hom}}\nolimits}
\def\Ext{\mathop{\mathrm{Ext}}\nolimits}
\def\SheafHom{\mathop{\mathcal{H}\!\mathit{om}}\nolimits}
\def\SheafExt{\mathop{\mathcal{E}\!\mathit{xt}}\nolimits}
\def\Sch{\mathit{Sch}}
\def\Mor{\mathop{\mathrm{Mor}}\nolimits}
\def\Ob{\mathop{\mathrm{Ob}}\nolimits}
\def\Sh{\mathop{\mathit{Sh}}\nolimits}
\def\NL{\mathop{N\!L}\nolimits}
\def\CH{\mathop{\mathrm{CH}}\nolimits}
\def\proetale{{pro\text{-}\acute{e}tale}}
\def\etale{{\acute{e}tale}}
\def\QCoh{\mathit{QCoh}}
\def\Ker{\mathop{\mathrm{Ker}}}
\def\Im{\mathop{\mathrm{Im}}}
\def\Coker{\mathop{\mathrm{Coker}}}
\def\Coim{\mathop{\mathrm{Coim}}}

% Boxtimes
%
\DeclareMathSymbol{\boxtimes}{\mathbin}{AMSa}{"02}

%
% Macros for moduli stacks/spaces
%
\def\QCohstack{\mathcal{QC}\!\mathit{oh}}
\def\Cohstack{\mathcal{C}\!\mathit{oh}}
\def\Spacesstack{\mathcal{S}\!\mathit{paces}}
\def\Quotfunctor{\mathrm{Quot}}
\def\Hilbfunctor{\mathrm{Hilb}}
\def\Curvesstack{\mathcal{C}\!\mathit{urves}}
\def\Polarizedstack{\mathcal{P}\!\mathit{olarized}}
\def\Complexesstack{\mathcal{C}\!\mathit{omplexes}}
% \Pic is the operator that assigns to X its picard group, usage \Pic(X)
% \Picardstack_{X/B} denotes the Picard stack of X over B
% \Picardfunctor_{X/B} denotes the Picard functor of X over B
\def\Pic{\mathop{\mathrm{Pic}}\nolimits}
\def\Picardstack{\mathcal{P}\!\mathit{ic}}
\def\Picardfunctor{\mathrm{Pic}}
\def\Deformationcategory{\mathcal{D}\!\mathit{ef}}
\usepackage{cleveref}
\crefname{originallemma}{original lemma}{original lemmas}
\Crefname{originallemma}{Original lemma}{Original lemmas}
\makeatletter
\newlabel{lemma-quotient-stack-functorial}{{046Q}{source}{}{sourcepointer-1}{}}
\newlabel{lemma-quotient-stack-functorial@cref}{{[originallemma][0][]046Q}{[1][0][]source}{}{}{}}
\makeatother
\crefname{originallemma}{original lemma}{original lemmas}
\Crefname{originallemma}{Original lemma}{Original lemmas}
\makeatletter
\newlabel{lemma-quotient-stack-arrows}{{044R}{source}{}{sourcepointer-2}{}}
\newlabel{lemma-quotient-stack-arrows@cref}{{[originallemma][0][]044R}{[1][0][]source}{}{}{}}
\makeatother
\crefname{originalreference}{original reference}{original references}
\Crefname{originalreference}{Original reference}{Original references}
\makeatletter
\newlabel{stacks-lemma-stackification-fibre-product-categories-fibred-in-groupoids}{{04Y2}{source}{}{sourcepointer-3}{}}
\newlabel{stacks-lemma-stackification-fibre-product-categories-fibred-in-groupoids@cref}{{[originalreference][0][]04Y2}{[1][0][]source}{}{}{}}
\makeatother
\title{The quotient-stack square of a groupoid morphism is a 2-fibre product}
\author{The Stacks Project Authors}
\date{stacks-project:spaces-groupoids.tex:local-snapshot-20261012; GNU FDL 1.2 or later}
\begin{document}
\maketitle
\noindent\textit{Editorial note.} All delivered mathematical text is the author's own, copied verbatim from the pinned Stacks Project chapter file spaces-groupoids.tex; no statement or proof step was rewritten, reordered or removed.\par\medskip
\noindent\textit{Editorial note.} Span recomputation by label pairing (LANE\_BRIEF 10.2): in this chapter the author's proof environment follows the statement environment as a standalone block; the spec covers the complete environment pair (statement span, proof span) as recorded in the pool candidate.\par\medskip
\ifdefined\hypertarget\hypertarget{sourcepointer-1}{}\fi
\noindent\textit{Source reference.} Original lemma 046Q: sources/stacks-chow/tags.txt, source lines 5481--5481. Cross-chapter standard prerequisite cited by the author inside the counted span; it is used as a pointer and its own proof is not part of this item.\par\medskip
\ifdefined\hypertarget\hypertarget{sourcepointer-2}{}\fi
\noindent\textit{Source reference.} Original lemma 044R: sources/stacks-chow/tags.txt, source lines 5410--5410. Cross-chapter standard prerequisite cited by the author inside the counted span; it is used as a pointer and its own proof is not part of this item.\par\medskip
\ifdefined\hypertarget\hypertarget{sourcepointer-3}{}\fi
\noindent\textit{Source reference.} Original reference 04Y2: sources/stacks-chow/tags.txt, source lines 6415--6415. Cross-chapter standard prerequisite cited by the author inside the counted span; it is used as a pointer and its own proof is not part of this item.\par\medskip
\begin{lemma}
\label{lemma-cartesian-square-of-morphism}
Notation and assumption as in
Lemma \ref{lemma-quotient-stack-functorial}.
Let $(U'', R'', s'', t'', c'')$ be the groupoid in algebraic spaces over $B$
constructed above.
There is a $2$-commutative square
$$
\xymatrix{
[U''/R''] \ar[d] \ar[r]_{[g]} & [U/R] \ar[d]^{[f]} \\
\mathcal{S}_{U'} \ar[r] & [U'/R']
}
$$
which identifies $[U''/R'']$ with the $2$-fibre product.
\end{lemma}

\begin{proof}
The maps $[f]$ and $[g]$ come from an application of
Lemma \ref{lemma-quotient-stack-functorial}
and the other two maps come from
Lemma \ref{lemma-quotient-stack-arrows}
(and the fact that $(U'', R'', s'', t'', c'')$ lives over $U'$).
To show the $2$-fibre product property, it suffices to prove the lemma
for the diagram
$$
\xymatrix{
[U''/_{\!p}R''] \ar[d] \ar[r]_{[g]} & [U/_{\!p}R] \ar[d]^{[f]} \\
\mathcal{S}_{U'} \ar[r] & [U'/_{\!p}R']
}
$$
of categories fibred in groupoids, see
Stacks, Lemma
\ref{stacks-lemma-stackification-fibre-product-categories-fibred-in-groupoids}.
In other words, it suffices to show that an object of the $2$-fibre product
$\mathcal{S}_U \times_{[U'/_{\!p}R']} [U/_{\!p}R]$ over $T$ corresponds
to a $T$-valued point of $U''$ and similarly for morphisms. And of course
this is exactly how we constructed $U''$ and $R''$ in the first place.

\medskip\noindent
In detail, an object of $\mathcal{S}_U \times_{[U'/_{\!p}R']} [U/_{\!p}R]$
over $T$ is a triple $(u', u, r')$ where $u'$ is a $T$-valued point of $U'$,
$u$ is a $T$-valued point of $U$, and $r'$ is a morphism from $u'$ to
$f(u)$ in $[U'/R']_T$, i.e., $r'$ is a $T$-valued point of $R$ with
$s'(r') = u'$ and $t'(r') = f(u)$. Clearly we can forget about $u'$
without losing information and we see that these objects are in one-to-one
correspondence with $T$-valued points of $R''$.

\medskip\noindent
Similarly for morphisms: Let $(u'_1, u_1, r'_1)$ and $(u'_2, u_2, r'_2)$ be
two objects of the fibre product over $T$. Then a morphism from
$(u'_2, u_2, r'_2)$ to $(u'_1, u_1, r'_1)$
is given by $(1, r)$ where $1 : u'_1 \to u'_2$ means simply $u'_1 = u'_2$
(this is so because $\mathcal{S}_U$ is fibred in sets), and $r$ is a
$T$-valued point of $R$ with $s(r) = u_2$, $t(r) = u_1$ and moreover
$c'(f(r), r'_2) = r'_1$. Hence the arrow
$$
(1, r) : (u'_2, u_2, r'_2) \to (u'_1, u_1, r'_1)
$$
is completely determined by knowing the pair $(r, r'_2)$. Thus the functor
of arrows is represented by $R''$, and moreover the morphisms $s''$, $t''$,
and $c''$ clearly correspond to source, target and composition in
the $2$-fibre product $\mathcal{S}_U \times_{[U'/_{\!p}R']} [U/_{\!p}R]$.
\end{proof}

\end{document}
